\section{Discussion}

\subsection{Main interpretation}
In this tournament, trailing states were conditionally associated with more frequent provider Pressure activity per opponent pass and higher estimated odds that a Pressure sequence was followed by controlled possession within five seconds. The two associations coexist, but neither identifies why the pattern appeared. Score state, opponent protection of a lead, field position, substitutions, and opportunity selection evolve together.

\subsection{Intensity and efficiency are distinct}
The simultaneous positive trailing contrasts are analytically meaningful because they arise from different numerators and denominators. Intensity asks how often Pressure events occurred given opponent passing opportunities. Efficiency asks how often a sequence converted given Pressure episodes. Tactical urgency, greater pressing commitment, altered field position, opponents accepting lower-risk possession, and selective initiation of Pressure are plausible hypotheses. They were not directly measured as mechanisms and should be tested with richer designs.

\subsection{Leading states}
Leading estimates were modestly negative relative to drawing for both primary outcomes, with intervals that included the reference-equivalent value on the coefficient scale. This does not establish equivalence or absence of a difference. It indicates that the available tournament sample produced less precise and smaller adjusted contrasts than for trailing.

\subsection{Measurement implications}
Pressure events, Pressure sequences, and classic PPDA answer different questions. Event counts reflect provider annotations; sequences reduce repeated actions into episodes; classic PPDA maps a documented defensive-action set to opponent passes and is inversely scaled. Denominator choice changes both reliability and interpretation. Fixed windows offer comparable time blocks but discard mixed states; homogeneous segments preserve state purity with varying duration. Analysts should state these choices before comparing teams or contexts.

\subsection{Practical relevance}
For analysts and coaches, the framework suggests monitoring activity, conversion, and post-regain value separately. A high activity rate need not imply efficient conversion, and efficient conversion need not imply valuable attacks. The estimates do not prescribe a tactical response. They provide benchmarks for asking whether an observed change reflects opportunity volume, episode conversion, or subsequent attack quality.

\subsection{Future 2026 extension}
A future milestone may connect pressing decisions with dynamic qualification pressure under the 48-team 2026 format. Such work would require prospective qualification-state modeling and safeguards against future-information leakage. No 2026 simulation or empirical result is included here.
